\section{Spatial Variation}\label{spatial}

The ecolab model can run in multicellular mode by assigning to the
{\tt dims}. The model in equation \ref{lotka-volterra} is replicated
in each cell. The idea behind this is to set up a spatially dependent
model. In case of two dimensional space, one can specify the x and y
dimensions by assigning the following values as part of the input
parameters \S\ref{input parameters}:
\begin{verbatim}
set dims(x) 10
set dims(y) 20
\end{verbatim}
This specifies a $10\times20$ grid of 200 cells.
The array indices ``x'' and ``y'' are purely arbitrary, one could use
``0'' and ``1'' instead.

In the multicellular case, the number of species in each cell is set
to be the value of the TCL variable {\tt nsp}. In order to specify
different numbers of species in each cell at initialisation, one can
assign to the {\tt nsp} array, in the following syntax (eg a 2D
system):
\begin{verbatim}
set nsp(0,0) 10
set nsp(0,1) 12
set nsp(1,0) 15
set nsp(1,1) 13
\end{verbatim}

One can also initialise other state variables with spatial dependence
by appending cell indices, e.g.
\begin{verbatim}
set density(list,0,0) 100
set density(list,0,1) 100
set density(list,1,0) 100
set density(list,1,1) 100
\end{verbatim}
In this case, one may need to write a short TCL loop to intialise a
particular pattern. This feature is not yet implemented as of version 3.2

Migration is performed using the {\tt migrate}\ref{migrate} operator.

\subsection{Parallel Execution}

A multicellular model can be run on an MPI parallel system by
compiling with the MPI flags turned on. Ecolab distributes the cells
across the processors in a block fashion. A TCL command can be
written that executes on all processors by placing the {\tt PARALLEL}
directive at the start of the command. Otherwise (by default) all code
executes just on processor 0. An example of a useful utility function is
{\tt get\_all\_globals} which assembles a complete representation of
the whole system on processor 0. This uses the {\tt PARALLEL give\_me}
command to obtain the slave processor state variables. There is a
corresponding {\tt take\_it} command to push state variable out to all
the slaves. These expect the master process to broadcast the data, so
these routines should not be called from TCL scripts, otherwise Ecolab
will hang.

Work on using the spatial Ecolab is reported in
\cite{Standish98c}. The model showed good scalability, and even
superlinear speedup in some circumstances.

The parallel coding practices employed in Ecolab are experimental, and
may well change in the next release.
